\documentclass[12pt]{article}

\usepackage{amsmath, amsfonts}
\usepackage[margin=1cm]{geometry}

\begin{document}

\noindent
{\large\bfseries Number Theory}
\hfill
znz
\quad
\today

\hrule
\medskip

\begin{enumerate}
  \item
        Let
        \[
        f(t)=t^3+t+1,
        \qquad
        g(t)=t^2+t+1.
        \]Compute explicitly in \(\mathbb F_2[t]\):
        \[
        f+g,
        \]\[
        fg,
        \]and
        \[
        f^2.
        \]Write the final answers with coefficients reduced in \(\mathbb F_2\).
  \item
        Let \(F\) be a field.
        Determine all units of
        \[
        F[t].
        \]Prove your claim.
        You may use
        \[
        \deg(fg)=\deg f+\deg g
        \]for nonzero polynomials.
  \item
        List all monic quadratic polynomials
        \[
        f\in\mathbb F_2[t].
        \]For each one, determine whether it is reducible or irreducible.
        You may use the root criterion for quadratics.
        Be explicit about which elements of \(\mathbb F_2\) must be checked.
  \item
        Let
        \[
        f,g\in F[t]
        \]be nonzero polynomials and suppose
        \[
        f\mid g.
        \]Prove that
        \[
        \deg f\le\deg g.
        \]Then determine what additional conclusion follows if
        \[
        \deg f=\deg g.
        \]Do not assume the second conclusion; derive it from the definition of divisibility and the degree formula.
  \item
        Consider the two statements:
        \[
        6=2\cdot3
        \]in \(\mathbb Z\), and
        \[
        t^2+t=t(t+1)
        \]in \(\mathbb F_2[t]\).
        Explain precisely, using the definitions of divisibility and irreducibility, in what sense
        \[
        t
        \quad\text{and}\quad
        t+1
        \]play roles analogous to prime factors.
        Then answer:
        Is \(t^2+t+1\) irreducible over \(\mathbb F_2\)?

        Give a proof, not merely a factor search.
\end{enumerate}

\end{document}
